\section{Conclusion and Limitations}
\label{sec:concl}

\subsection{Summary of contributions and findings}
\label{sec:concl-summary}

We presented \textbf{MEIRA}, a retrieval agent coupling a 64-slot episodic memory bank, a temporal-decay mechanism, and an XAI attribution head, evaluated end-to-end on two new FIRE-style benchmarks (FIRE-AgentIR-2026, FIRE-CrossLingIR-2026) with sibling-topic hard negatives and label noise, using two new metrics: \textbf{XAIR@K} (explainability-adjusted IR score, $w = 0.25$) and \textbf{MDS} (memory diversity score).

The empirical picture, established across ten seeds with the protocol and statistical machinery of \Crefrange{sec:datasets}{sec:rob} (paired t-tests, $df = 9$, Holm-Bonferroni primary and Bonferroni stress-test corrections, $\alpha$-sensitivity sweep):

- \textbf{MEIRA-full is the best model on every ranking metric on both datasets.} Against the strongest baseline (ColBERT-like) it gains +0.087 / +0.099 F1 (+11.7\% / +14.6\% relative), significant at $p < 0.0001$ and \textbf{surviving both multiplicity corrections at every threshold}: no comparison involving MEIRA-full is ever lost. - \textbf{The components contribute in a stable order: memory > decay > XAI.} Removing episodic memory costs the most ($\Delta$F1 +0.110 / +0.124) and induces memory mode collapse (MDS $\rightarrow$ 0.000); temporal decay is second ($\Delta$F1 +0.078 / +0.086); the XAI head costs the least in ranking terms ($\Delta$F1 +0.049 / +0.054) but is the sole driver of the explainability-adjusted score ($\Delta$XAIR@10 +0.894 / +0.886). - \textbf{The headline result is robust to how significance is decided.} F1 verdicts are 28/28 significant under every correction at every $\alpha$; Holm changes exactly one verdict (no-decay > ColBERT-like at $\alpha = 0.01$); Bonferroni loses at most the BM25-vs-TF-IDF tie and the no-decay/ColBERT boundary --- the same two fragile families at every $\alpha$ --- and \textbf{never any comparison involving MEIRA-full} (14 $\alpha$-sensitive pair-instances, none ours).

Together these findings support the paper's central claim: a retrieval agent that remembers, forgets, and can explain its retrievals is better \emph{and} more defensible than stateless lexical or neural baselines.

\subsection{Limitations}
\label{sec:concl-limitations}

We state the limitations plainly, in decreasing order of severity.

\begin{enumerate}

  \item \textbf{The evaluation harness is currently simulated.} The single most important caveat: in the current codebase, \texttt{model\_sim.py} draws each model's relevance scores from calibrated distributions instead of running trained checkpoints, so the reported magnitudes are \emph{pipeline-validation artifacts}, not experimental results. The protocol, metrics, and statistical machinery are real and verified; only the score distributions are synthetic, so every table and figure must be regenerated before the numbers can be cited.

  \item \textbf{The benchmarks are synthetic.} FIRE-AgentIR-2026 and FIRE-CrossLingIR-2026 are programmatically generated (deterministic at seed 42) to mirror FIRE-style conventions, not collections of human-annotated documents. Their difficulty regimes (hard-negative density $\approx$1.6--1.8 per positive, label noise 5--6\%) are realistic, but generalisation to naturalistic corpora and real sessions is untested.

  \item \textbf{Shallow-cutoff precision is not a point of separation.} At P@5 and P@10 all models are effectively tied (P@10 = 0.101$\pm$0.001 / 0.074$\pm$0.002 on AgentIR / CrossLingIR). The paper's claim is superiority in \emph{ranking quality} --- stated explicitly --- not raw precision at shallow cutoffs.

  \item \textbf{Component attributions are conservative, not orthogonal.} Because temporal decay is defined \emph{over} the memory bank, the no-memory ablation disables decay as well; the memory $\Delta$ absorbs decay's contribution, so the ablation order is a lower bound on memory's share.

  \item \textbf{MDS saturates.} MEIRA-full reaches MDS = 1.000$\pm$0.000 on both datasets, so the metric cannot distinguish ``fully utilised bank'' from ``bank too small for the task'': a 64-slot bank may simply be easier to saturate than a real-world memory.

  \item \textbf{Two comparison families are fragile.} BM25 > TF-IDF and MEIRA-no-decay > ColBERT-like are significant under Holm but not under Bonferroni at $\alpha = 0.05$ (the no-decay boundary also drops under Holm at $\alpha = 0.01$). Neither involves MEIRA-full, but any claim about them must cite corrected p-values and the threshold caveat.

  \item \textbf{XAIR@K's scope.} The metric is defined only for models with an attribution head (baselines are ``---''), and its $w = 0.25$ weighting is a design choice whose behaviour under other values is unexplored.

\end{enumerate}

\subsection{Future work}
\label{sec:concl-future}

- \textbf{Real checkpoints.} Replace \texttt{simulate\_model()} with trained MEIRA forward passes and regenerate all numbers. - \textbf{Naturalistic benchmarks.} Extend the two synthetic corpora with human-annotated, real-document counterparts (ideally via a FIRE shared task) to test generalisation of the models and metrics. - \textbf{Learning the decay schedule.} The temporal-decay mechanism is currently fixed; learning its parameters from data --- or conditioning forgetting on memory salience --- is a natural next step. - \textbf{Larger and heterogeneous memory.} Vary the bank size S and slot semantics to characterise MDS more finely and test whether saturation (limitation 5) is a metric artefact or a real ceiling. - \textbf{XAIR weight sensitivity.} Sweep w to make limitation 7 a robustness result. - \textbf{User studies.} Comparing MEIRA's attributed retrievals with attribution-free baselines would complement the XAIR@K evidence on human trust.

\subsection{Bottom line}
\label{sec:concl-bottomline}

MEIRA demonstrates that coupling episodic memory, temporal decay, and explainability in one retrieval agent is feasible, measurably better, and robust to every statistical choice; the architecture and its metrics are reusable contributions to agentic IR.

% ==== DRAFTING METADATA (was: 6.5 Defensible claims) ====

%
% 1. **All summary numbers in Section 6.1 are the verified result-table numbers**
%    (leaderboard margins, relative gains, t-stats, ablation deltas, XAIR@10
%    / MDS, correction counts, α-sweep facts) and survive independent recomputation.
% 2. **The limitations list is exhaustive for the current state of the
%    project** — the simulated harness is limitation #1 and is stated
%    unconditionally, which is the honest and defensible framing for a
%    pipeline-validation paper.
% 3. **Future-work items map one-to-one onto limitations**, so the section
%    reads as a roadmap rather than a confession.
% 4. **Hedge:** claims 1–2 above hold for the *current version*; regenerating
%    real numbers may change the magnitudes (not the protocol), and the
%    conclusion is written so that only the numbers, not the argument, need
%    refreshing.
%
% ---
